% GENERATED FILE. Edit canonical lessons, then run npm run generate:books.
% canonical-source-hash: fnv1a64:38513e169612e2e8
% canonical-lessons: TE-C318-anipincu, TE-C318-anukuntunnanu, TE-C318-kabatti, TE-R318-opinion-and-result

\chapter{What I Think, and So}
\label{ch:te-what-i-think-and-so}
\hlchaptermodality{318}

\begin{chapteropening}
\textbf{By the end of this chapter:} \emph{I can say what I think, firmly or softly, with ani anukuṇṭunnānu, nā abhiprāyaṁlō and anipistundi, and give a result with kābaṭṭi and andukē.}

Complete TE-R318-opinion-and-result: give an opinion three ways and draw a result from a reason with kābaṭṭi and andukē.
\end{chapteropening}

\section[\texorpdfstring{\te{అనిపించు} (anipiñcu)}{అనిపించు (anipiñcu)}]{\te{అనిపించు} (anipiñcu) — to seem, to feel}

\label{lesson:TE-C318-anipincu}



\begin{quote}
Before the new one: say \textquotedblleft{}I came home and had a meal\textquotedblright{}. (\emph{nēnu iṇṭiki vacci bhōjanaṁ cēśānu}.) Then say what \textbf{\te{ఈ}} means on its own, in front of a noun. (\textquotedblleft{}This.\textquotedblright{})
\end{quote}

\subsection*{You'll want to know: \te{అనిపించు}}
\textbf{\te{అనిపించు}} — \emph{anipiñcu} — \textquotedblleft{}to seem, to feel (to someone)\textquotedblright{}.

\begin{quote}
\textbf{\te{మొదట} \te{తెలుగు} \te{నాకు} \te{కష్టంగా} \te{అనిపించింది}.}
\end{quote}

\begin{quote}
\emph{modaṭa telugu nāku kaṣṭaṅgā anipiñcindi.} — \textquotedblleft{}At first Telugu seemed hard to me.\textquotedblright{}
\end{quote}

\begin{quote}
\textbf{\te{ఇప్పుడు} \te{అది} \te{నాకు} \te{సులభంగా} \te{అనిపిస్తుంది}.}
\end{quote}

\begin{quote}
\emph{ippuḍu adi nāku sulabhaṅgā anipistundi.} — \textquotedblleft{}Now it seems easy to me.\textquotedblright{}
\end{quote}

The one who feels sits in the dative, \textbf{\te{నాకు}}, the way it does in \textbf{\te{నాకు} \te{తెలుగు} \te{వచ్చు}}. The verb stays \textquotedblleft{}it\textquotedblright{}: \textbf{\te{అనిపిస్తుంది}} (\emph{anipistundi}, \textquotedblleft{}it seems\textquotedblright{}), \textbf{\te{అనిపించింది}} (\emph{anipiñcindi}, \textquotedblleft{}it seemed\textquotedblright{}). The quality takes \textbf{-\te{గా}}: \textbf{\te{కష్టంగా}}, \textbf{\te{సులభంగా}}. \textbf{\te{నాకు} \te{అలా} \te{అనిపిస్తుంది}} is \textquotedblleft{}that is how it seems to me\textquotedblright{}: \textbf{\te{అలా}} (\emph{alā}, \textquotedblleft{}like that\textquotedblright{}) is the \textbf{\te{అలా}} inside \textbf{\te{అలాగే}}.

\begin{cousinweb}
\textbf{\te{అనిపించు}} (\emph{anipiñcu}) is native Dravidian. DEDR 868 lists Telugu \emph{anu}, \textquotedblleft{}to say\textquotedblright{}, and beside it \emph{anipincu}, \textquotedblleft{}to seem, appear, be considered\textquotedblright{}, which it also glosses \textquotedblleft{}to cause to be said\textquotedblright{}. So \textbf{\te{నాకు} \te{అలా} \te{అనిపిస్తుంది}} is close to \textquotedblleft{}it gets said to me like that\textquotedblright{}. The same \textbf{\te{అను}} is inside \textbf{\te{అనుకో}}, \textquotedblleft{}to think\textquotedblright{}, and \textbf{\te{ఎందుకంటే}}.
\end{cousinweb}

\subsection*{Your turn}
Say these aloud:

\begin{itemize}
\raggedright
  \item \emph{anipiñcu}, slowly, then at speed
  \item \textquotedblleft{}At first Telugu seemed hard to me\textquotedblright{} — \emph{modaṭa telugu nāku kaṣṭaṅgā anipiñcindi}
  \item \textquotedblleft{}Now it seems easy to me\textquotedblright{} — \emph{ippuḍu adi nāku sulabhaṅgā anipistundi}
  \item \textquotedblleft{}That is how it seems to me\textquotedblright{} — \emph{nāku alā anipistundi}
\end{itemize}

\begin{tcolorbox}[breakable,colback=teal!4,colframe=teal!35!black,title={Before you move on}]
What does \textbf{\te{అనిపించు}} mean? (\textbf{To seem, to feel.}) Where does the one who feels go? (In the dative: \textbf{\te{నాకు}}.) \textquotedblleft{}It seemed\textquotedblright{}? (\textbf{\te{అనిపించింది}}.)
\end{tcolorbox}
\section[\texorpdfstring{\te{అని} \te{అనుకుంటున్నాను} (ani anukuṇṭunnānu)}{అని అనుకుంటున్నాను (ani anukuṇṭunnānu)}]{\te{అని} \te{అనుకుంటున్నాను} (ani anukuṇṭunnānu) — I think that}

\label{lesson:TE-C318-anukuntunnanu}



\begin{quote}
Say \textquotedblleft{}to think\textquotedblright{} (\emph{anukō}) and \textquotedblleft{}an opinion\textquotedblright{} (\emph{abhiprāyaṁ}). Then \textquotedblleft{}it seems so to me\textquotedblright{}. (\emph{nāku alā anipistundi}.)
\end{quote}

\subsection*{You'll want to know: … \te{అని} \te{అనుకుంటున్నాను}}
\begin{quote}
\textbf{\te{తెలుగు} \te{కష్టం} \te{కాదని} \te{నేను} \te{అనుకుంటున్నాను}.}
\end{quote}

\begin{quote}
\emph{telugu kaṣṭaṁ kādani nēnu anukuṇṭunnānu.} — \textquotedblleft{}I think Telugu is not hard.\textquotedblright{}
\end{quote}

The thought comes first, closed by \textbf{\te{అని}} (\emph{ani}, \textquotedblleft{}saying\textquotedblright{}); then the verb of thinking. A final \textbf{-\te{ఉ}} melts into \textbf{\te{అని}}: \textbf{\te{కాదు}} + \textbf{\te{అని}} $\to$ \textbf{\te{కాదని}}; \textbf{\te{పడుతుంది}} + \textbf{\te{అని}} $\to$ \textbf{\te{పడుతుందని}}.

\begin{quote}
\textbf{\te{రేపు} \te{వర్షం} \te{పడుతుందని} \te{నేను} \te{అనుకుంటున్నాను}.}
\end{quote}

\begin{quote}
\emph{rēpu varṣaṁ paḍutundani nēnu anukuṇṭunnānu.} — \textquotedblleft{}I think it will rain tomorrow.\textquotedblright{}
\end{quote}

\textbf{\te{అనుకుంటున్నాను}} is the \textbf{-\te{తున్నాను}} form of \textbf{\te{అనుకో}}: \textquotedblleft{}I am thinking (that)\textquotedblright{}, the everyday \textquotedblleft{}I think\textquotedblright{}.

\subsection*{You'll want to know: \te{నా} \te{అభిప్రాయంలో}}
\begin{quote}
\textbf{\te{నా} \te{అభిప్రాయంలో} \te{ఈ} \te{హోటలు} \te{చాలా} \te{మంచిది}.}
\end{quote}

\begin{quote}
\emph{nā abhiprāyaṁlō ī hōṭalu cālā mañcidi.} — \textquotedblleft{}In my opinion this restaurant is very good.\textquotedblright{}
\end{quote}

\begin{quote}
\textbf{\te{నాకు} \te{అలా} \te{అనిపిస్తుంది}.}
\end{quote}

\begin{quote}
\emph{nāku alā anipistundi.} — \textquotedblleft{}That is how it seems to me.\textquotedblright{}
\end{quote}

Three ways, from firm to soft. \textbf{\te{నా} \te{అభిప్రాయంలో}} is careful and a little formal. \textbf{… \te{అని} \te{అనుకుంటున్నాను}} is the plain one. \textbf{\te{నాకు} … \te{అనిపిస్తుంది}} is the softest: a feeling more than a claim. \textbf{\te{నా} \te{ఉద్దేశంలో}}, with the \textbf{\te{ఉద్దేశం}} you know, is heard too.

\subsection*{Your turn}
Say these aloud:

\begin{itemize}
\raggedright
  \item \textquotedblleft{}I think Telugu is not hard\textquotedblright{} — \emph{telugu kaṣṭaṁ kādani nēnu anukuṇṭunnānu}
  \item \textquotedblleft{}I think it will rain tomorrow\textquotedblright{} — \emph{rēpu varṣaṁ paḍutundani nēnu anukuṇṭunnānu}
  \item \textquotedblleft{}In my opinion this restaurant is very good\textquotedblright{} — \emph{nā abhiprāyaṁlō ī hōṭalu cālā mañcidi}
  \item your own opinion of a film you saw, three ways: firm, plain and soft
\end{itemize}

\begin{tcolorbox}[breakable,colback=teal!4,colframe=teal!35!black,title={Before you move on}]
Where does the thought go, before or after \textbf{\te{అని}}? (\textbf{Before.}) \textquotedblleft{}Is not\textquotedblright{} + \textbf{\te{అని}}? (\textbf{\te{కాదని}}.) \textquotedblleft{}In my opinion\textquotedblright{}? (\textbf{\te{నా} \te{అభిప్రాయంలో}}.) Which way is the softest? (\textbf{\te{నాకు} … \te{అనిపిస్తుంది}}.)
\end{tcolorbox}
\section[\texorpdfstring{\te{కాబట్టి} (kābaṭṭi)}{కాబట్టి (kābaṭṭi)}]{\te{కాబట్టి} (kābaṭṭi) — so: the reason, then the result}

\label{lesson:TE-C318-kabatti}



\begin{quote}
Say \textquotedblleft{}therefore\textquotedblright{} (\emph{anduvalla}), \textquotedblleft{}that is exactly why\textquotedblright{} (\emph{andukē}) and \textquotedblleft{}because\textquotedblright{} (\emph{endukaṇṭē}).
\end{quote}

\subsection*{You'll want to know: a reason + \te{కాబట్టి}}
\begin{quote}
\textbf{\te{వర్షం} \te{పడుతోంది} \te{కాబట్టి} \te{నేను} \te{ఇంట్లో} \te{ఉన్నాను}.}
\end{quote}

\begin{quote}
\emph{varṣaṁ paḍutōndi kābaṭṭi nēnu iṇṭlō unnānu.} — \textquotedblleft{}It is raining, so I am at home.\textquotedblright{}
\end{quote}

\begin{quote}
\textbf{\te{బస్సు} \te{లేదు} \te{కాబట్టి} \te{మేము} \te{నడిచాము}.}
\end{quote}

\begin{quote}
\emph{bassu lēdu kābaṭṭi mēmu naḍicāmu.} — \textquotedblleft{}There was no bus, so we walked.\textquotedblright{}
\end{quote}

\textbf{\te{కాబట్టి}} follows the reason directly, inside the same sentence; the result comes after it. Read it as \textquotedblleft{}since it is raining, …\textquotedblright{} or \textquotedblleft{}it is raining, so …\textquotedblright{}. The order is the reverse of \textbf{\te{ఎందుకంటే}}, where the result comes first and the reason follows.

\subsection*{You'll want to know: \te{అందుకే}, a new sentence}
\begin{quote}
\textbf{\te{నాకు} \te{జ్వరం} \te{వచ్చింది}. \te{అందుకే} \te{నేను} \te{బడికి} \te{వెళ్ళలేదు}.}
\end{quote}

\begin{quote}
\emph{nāku jvaraṁ vaccindi. andukē nēnu baḍiki veḷḷalēdu.} — \textquotedblleft{}I had a fever. That is why I did not go to school.\textquotedblright{}
\end{quote}

When the reason is a whole sentence, the next sentence opens on \textbf{\te{అందుకే}}, \textquotedblleft{}that is exactly why\textquotedblright{}, or, more formally, \textbf{\te{అందువల్ల}}, \textquotedblleft{}therefore\textquotedblright{}.

\textbf{\te{వెళ్ళలేదు}} (\emph{veḷḷalēdu}, \textquotedblleft{}did not go\textquotedblright{}) is how a story says \textquotedblleft{}did not\textquotedblright{}: the verb's base with \textbf{\te{లేదు}}. \textbf{\te{చేయలేదు}}, \textquotedblleft{}did not do\textquotedblright{}; \textbf{\te{రాలేదు}}, \textquotedblleft{}did not come\textquotedblright{}.

\subsection*{Your turn}
Say these aloud:

\begin{itemize}
\raggedright
  \item \textquotedblleft{}It is raining, so I am at home\textquotedblright{} — \emph{varṣaṁ paḍutōndi kābaṭṭi nēnu iṇṭlō unnānu}
  \item \textquotedblleft{}There was no bus, so we walked\textquotedblright{} — \emph{bassu lēdu kābaṭṭi mēmu naḍicāmu}
  \item \textquotedblleft{}I had a fever. That is why I did not go to school.\textquotedblright{} — \emph{nāku jvaraṁ vaccindi. andukē nēnu baḍiki veḷḷalēdu.}
  \item one thing you did not do yesterday, and why, with \te{కాబట్టి}
\end{itemize}

\begin{tcolorbox}[breakable,colback=teal!4,colframe=teal!35!black,title={Before you move on}]
With \textbf{\te{కాబట్టి}}, which comes first, the reason or the result? (\textbf{The reason.}) With \textbf{\te{ఎందుకంటే}}? (\textbf{The result.}) \textquotedblleft{}Did not come\textquotedblright{}? (\textbf{\te{రాలేదు}}.)
\end{tcolorbox}
\section[(review)]{Review — opinion and result}

\label{lesson:TE-R318-opinion-and-result}



\begin{quote}
Say \textquotedblleft{}it seemed hard to me\textquotedblright{}. (\emph{nāku kaṣṭaṅgā anipiñcindi}.) Then say \textquotedblleft{}in the meantime\textquotedblright{}. (\emph{īlōgā}.)

\begin{itemize}
\raggedright
  \item \emph{Write it:} \textbf{\te{ఈ}}, the first letter of \textbf{\te{ఈలోగా}}, once
\end{itemize}
\end{quote}

\subsection*{Your turn}
Say these aloud:

\begin{itemize}
\raggedright
  \item I think this book is good.
  \item In my opinion Telugu is not hard.
  \item It is raining, so we did not go to the town.
  \item I had a fever. That is why I stayed at home.
  \item First we had a meal. After that we went to a film. In the end we came home.
  \item the reading's trip in two sentences, with -\te{ఇ} forms
\end{itemize}

\begin{tcolorbox}[breakable,colback=teal!4,colframe=teal!35!black,title={Before you move on}]
\textquotedblleft{}I think that …\textquotedblright{}? (\textbf{… \te{అని} \te{అనుకుంటున్నాను}}.) \textquotedblleft{}In my opinion\textquotedblright{}? (\textbf{\te{నా} \te{అభిప్రాయంలో}}.) \textquotedblleft{}It seems to me\textquotedblright{}? (\textbf{\te{నాకు} … \te{అనిపిస్తుంది}}.) \textquotedblleft{}So\textquotedblright{}, after a reason? (\textbf{\te{కాబట్టి}}.) \textquotedblleft{}That is why\textquotedblright{}, opening a sentence? (\textbf{\te{అందుకే}}.)
\end{tcolorbox}
